\section*{Bab \ref{ch:b1:angular-momentum} --- Momentum Sudut}

\begin{solution}{exo:b1:angular-momentum:1}
$Fd = 50 \times 0.30 = \qty{15}{N.m}$; $Fd\sin 60^\circ = \qty{13}{N.m}$;
menyusuri gagangnya, garis kerjanya melalui murnya: $0$.
\end{solution}

\begin{solution}{exo:b1:angular-momentum:2}
$L = mrv = 5.97\times10^{24} \times 1.50\times10^{11} \times 2.98\times10^4 =
\qty{2.7e40}{kg.m^2/s}$; $\dd A/\dd t = L/2m = rv/2 = \qty{2.2e15}{m^2/s}$.
\end{solution}

\begin{solution}{exo:b1:angular-momentum:3}
$r = \ell\sin 30^\circ = \qty{0.50}{m}$, $\omega = \qty{3.4}{rad/s}$:
$L_z/m = r^2\omega = \qty{0.84}{m^2/s}$. Garis kerja tegangannya melalui
porosnya (tak bermomen terhadap sumbu mana pun yang melaluinya);
beratnya sejajar sumbu tegaknya (tak bermomen terhadapnya): $L_z$ kekal.
\end{solution}

\begin{solution}{exo:b1:angular-momentum:4}
$J_1 = 4.0 + 2 \times 3.0 \times 0.64 = \qty{7.8}{kg.m^2}$, $J_2 = 4.0 +
2 \times 3.0 \times 0.04 = \qty{4.2}{kg.m^2}$; $\omega_2 = \omega_1J_1/J_2 =
1.85\omega_1$: \qty{3.7}{} putaran tiap sekon. $E_k = L^2/2J$: nisbahnya
$J_1/J_2 = 1.85$; ototnya, yang menarik lengannya ke dalam melawan
kecenderungan sentrifugalnya, yang melakukan usahanya.
\end{solution}

\begin{solution}{exo:b1:angular-momentum:5}
$L_\Delta = m\ell^2\dot\theta$; tegangannya: tak bermomen (memotong
sumbunya); beratnya: $-mg\ell\sin\theta$ (\omterm{def:b1:angular-momentum:moment}{lengan gayanya}
$\ell\sin\theta$, memulihkan). $m\ell^2\ddot\theta = -mg\ell\sin\theta$.
\end{solution}

\begin{solution}{exo:b1:angular-momentum:6}
$r_pv_p = r_av_a$: $v_p/v_a = r_a/r_p = 1.034$; dengan $(v_p + v_a)/2 \approx
\qty{29.8}{km/s}$: $v_p = \qty{30.3}{km/s}$, $v_a = \qty{29.3}{km/s}$.
\end{solution}

\begin{solution}{exo:b1:angular-momentum:7}
$L = mrv$ kekal: $v_2 = v_1r_1/r_2 = \qty{4.0}{m/s}$, $\omega_2 = v_2/r_2
= \qty{16}{rad/s}$ (dari \qty{4}{rad/s}). $E_1 = \qty{0.40}{J}$, $E_2 =
\qty{1.6}{J}$: tarikannya melakukan usaha \qty{1.2}{J}. Tegangannya
memang radial, tetapi perpindahan kepingnya kini mempunyai komponen
radial (ia memilin ke dalam): usaha $\vect T\cdot\dd\vect\ell$ tak lagi nol.
\end{solution}

\begin{solution}{exo:b1:angular-momentum:8}
\omterm{def:b1:angular-momentum:moment}{Momen} terhadap tumpuannya (panjangnya dari tumpuannya): anaknya $30 \times 1.5 =
45$ (kiri); berat papannya di pusatnya, $0.5$ ke kanan: $20 \times
0.5 = 10$ (kanan); anak keduanya pada $x$ ke kanan: $45x$. Setimbang:
$45 = 10 + 45x$, $x = \qty{0.78}{m}$. Tumpuannya menopang berat totalnya,
$95 \times 9.81 = \qty{932}{N}$.
\end{solution}

\begin{solution}{exo:b1:angular-momentum:9}
Massanya: $m_1a = m_1g - T_1$, $m_2a = T_2 - m_2g$; katrolnya: $J\dot\omega =
(T_1 - T_2)R$ dengan $a = R\dot\omega$. Dijumlahkan: $a = (m_1 - m_2)g/(m_1 + m_2 +
J/R^2) = 9.81/(3.0 + 1.0) = \qty{2.5}{m/s^2}$ (terhadap \qty{3.3}{m/s^2});
$T_1 = m_1(g - a) = \qty{14.7}{N}$, $T_2 = m_2(g + a) = \qty{12.3}{N}$ ---
tak sama besar, karena katrolnya memerlukan momen total.
\end{solution}

\begin{solution}{exo:b1:angular-momentum:10}
$r_1v_1 = r_2v_{\perp 2}$: $v_{\perp 2} = 40 \times 0.5/10 = \qty{2.0}{km/s}$.
\omterm{cor:b1:angular-momentum:central}{Hukum luas sapuan} hanya memberi komponen tegak lurusnya; komponen
radialnya menuntut kekekalan energi (\cref{ch:b1:central-forces}).
\end{solution}

\begin{solution}{exo:b1:angular-momentum:11}
$J\omega$ tetap dengan $J \propto R^2$: $T' = T(R'/R)^2 = 25 \times 86400
\times (10^4/7\times10^8)^2 = \qty{4.4e-4}{s}$. Laju khatulistiwanya $2\pi R'/T'
= \qty{1.4e8}{m/s}$, separuh laju cahaya: pengerutannya tak dapat menjaga
seluruh \omterm{def:b1:angular-momentum:L}{momentum sudutnya} --- banyak yang dilepaskan (dan kerelatifan masuk).
\end{solution}

\begin{solution}{exo:b1:angular-momentum:12}
$L_z = mr^2\omega$, dan $r$ berubah: $\dd L_z/\dd t = 2mr\dot r\omega \neq 0$.
Ini pastilah momen reaksi melintang batangnya $N$: $rN =
2mr\dot r\omega$, sehingga $N = 2m\dot r\omega$ --- batangnya mendorong
maniknya ke samping selagi ia meluncur ke luar (suku Coriolis
\cref{ch:b1:non-inertial-frames}).
\end{solution}

\begin{solution}{pb:b1:angular-momentum:1}
\textbf{1.} $m\omega^2r = GMm/r^2$: $\omega = \sqrt{GM/r^3} = \sqrt{3.98\times
10^{14}/5.66\times10^{25}} = \qty{2.65e-6}{rad/s}$, periodenya $2\pi/\omega =
\qty{2.37e6}{s} = \qty{27.4}{d}$.

\textbf{2.} $L_{\text{orb}} = 7.35\times10^{22} \times 1.475\times10^{17}
\times 2.65\times10^{-6} = \qty{2.9e34}{kg.m^2/s}$.

\textbf{3.} $v = \sqrt{GM/r}$, $L = mrv = m\sqrt{GMr}$; $\dd L/\dd r =
\tfrac12 m\sqrt{GM/r} = L/2r$.

\textbf{4.} $J = 0.33 \times 5.97\times10^{24} \times 4.06\times10^{13} =
\qty{8.0e37}{kg.m^2}$; $\Omega = \qty{7.29e-5}{rad/s}$; $L_{\text{spin}} =
\qty{5.8e33}{kg.m^2/s}$, seperlima $L_{\text{orb}}$.

\textbf{5.} $J_m \approx 0.4 \times 7.35\times10^{22} \times 3.0\times10^{12} =
\qty{8.9e34}{kg.m^2}$, dikali $\omega$: \qty{2.4e29}{kg.m^2/s}, yaitu $10^{-5}$
dari milik orbitnya: dapat diabaikan.

\textbf{6.} Tarikan Matahari bekerja pada pusat massa sistemnya (tak
bermomen terhadapnya sampai orde pertama); pengaruh pasang surutnya pada
pasangan itu kecil: total $L$ kekal sampai ketelitian yang diperlukan di sini.

\textbf{7.} Tonjolan yang lebih dekat, yang mendahului garis Bumi--Bulan,
ditarik Bulan dengan gaya yang komponen tangensialnya melawan putaran Bumi.

\textbf{8.} \omterm{def:b1:angular-momentum:moment}{Momen} negatif pada putarannya: Bumi melambat, harinya memanjang.

\textbf{9.} Tonjolannya menarik Bulan maju: momen positif pada orbitnya,
$L_{\text{orb}}$ tumbuh.

\textbf{10.} $L_{\text{orb}} \propto \sqrt r$ tumbuh, sehingga $r$ tumbuh; $v \propto
1/\sqrt r$ menurun.

\textbf{11.} Tarikan majunya melakukan usaha positif, yang menaikkan
energi Bulan; orbit yang lebih tinggi adalah orbit yang lebih lambat:
energinya pergi menjadi ketinggian (potensial), lebih banyak daripada
\omterm{thm:b1:work-and-energy:ke}{energi kinetik} yang hilang.

\textbf{12.} Penguncian pasang surut Bumi: harinya sama dengan bulannya,
tak ada tonjolan yang mendahului, tak ada \omterm{def:b1:angular-momentum:moment}{torsi}.

\textbf{13.} $\Delta\Omega/\Omega = -\qty{2.3e-3}{s}/\qty{86164}{s} =
\num{-2.7e-8}$ tiap abad.

\textbf{14.} $\Delta L_{\text{spin}} = L_{\text{spin}}\Delta\Omega/\Omega =
\qty{-1.55e26}{kg.m^2/s}$ tiap abad.

\textbf{15.} $\Delta L_{\text{orb}} = +\qty{1.55e26}{kg.m^2/s}$ tiap abad.

\textbf{16.} $\Delta r = 2r\,\Delta L/L = 2 \times 3.84\times10^8 \times 1.55
\times10^{26}/2.9\times10^{34} = \qty{4.1}{m}$ tiap abad, \qty{4.1}{cm/yr}:
dekat dengan \qty{3.8}{cm/yr} yang terukur (memanjangnya hari juga
mempunyai bagian yang bukan pasang surut).

\textbf{17.} $\Delta T/T = \tfrac32\Delta r/r = 1.6\times10^{-8}$: $2.37\times
10^6 \times 1.6\times10^{-8} = \qty{0.04}{s}$ tiap abad.

\textbf{18.} $\Delta E_{\text{rot}} = L_{\text{spin}}\Delta\Omega = 5.8\times10^{33}
\times 7.29\times10^{-5} \times (-2.7\times10^{-8}) = \qty{-1.1e22}{J}$ tiap
abad.

\textbf{19.} $\Delta E_{\text{orb}} = GMm\Delta r/2r^2 = 2.93\times10^{37}
\times 4.1/(2 \times 1.475\times10^{17}) = \qty{+4.1e20}{J}$ tiap abad.

\textbf{20.} Bersihnya $-1.1\times10^{22} + 4\times10^{20} \approx \qty{-1.06e22}{J}$
tiap abad: $3.4\times10^{12}\ \unit{W}$, yang dilesapkan sebagai kalor
oleh gesekan pasang surut di lautan dan di Bumi padatnya.

\textbf{21.} Seperlima daya umat manusia, seperempat belas aliran panas
buminya --- pemanas yang kecil tetapi tetap ada.

\textbf{22.} $L_{\text{tot}} = 3.5\times10^{34}$. Pada $r = 5.5\times10^8$:
$\omega = \sqrt{GM/r^3} = \qty{1.55e-6}{rad/s}$, $J\omega = 1.2\times10^{32}$,
$m\sqrt{GMr} = 7.35\times10^{22} \times 4.68\times10^{11} = 3.44\times10^{34}$;
jumlahnya $3.45\times10^{34}$: taat asas. Periodenya $2\pi/\omega = \qty{4.1e6}{s}
= \qty{47}{d}$.

\textbf{23.} Pasang surut Matahari terus memperlambat Bumi sampai di
bawah laju orbit Bulan, sehingga \omterm{def:b1:angular-momentum:moment}{torsi} Bumi--Bulan lalu akan berbalik;
lagi pula Matahari akan sudah berevolusi jauh sebelum itu (skala
waktunya puluhan miliar tahun).

\textbf{24.} Bumi membangkitkan pasang surut di Bulan, yang mengerem
putarannya sampai satu wajahnya tetap menghadap kita; karena lebih
ringan dan lebih dekat dengan pasangan yang lebih berat, penguncian
Bulan jauh lebih cepat.

\textbf{25.} Yang kekal: \omterm{def:b1:angular-momentum:L}{momentum sudut} totalnya. Yang dipertukarkan:
\omterm{def:b1:angular-momentum:L}{momentum sudut} dari putaran Bumi ke orbit Bulan, sekitar \qty{1.6e26}{SI}
tiap abad. Yang dilesapkan: \omterm{thm:b1:work-and-energy:em}{energi mekanik}, beberapa terawatt, sebagai
kalor pasang surut.
\end{solution}
